\Askhsh{1811}{4}
\begin{ylh}{1}
    \item 3.2. 1ο Κριτήριο ισότητας τριγώνων
    \item 4.2. Τέμνουσα δύο ευθειών - Ευκλείδειο αίτημα
    \item 4.6. Άθροισμα γωνιών τριγώνου
    \item 5.9. Μια ιδιότητα του ορθογώνιου τριγώνου.
\end{ylh}
Δίνονται δυο παράλληλες ευθείες $(\varepsilon)$ και $(\eta)$, και μια τρίτη που τις τέμνει στα σημεία $\mathrm{A}$ και $\mathrm{B}$ αντίστοιχα. Θεωρούμε τις διχοτόμους των εντός και επί τα αυτά μέρη γωνιών που σχηματίζονται, οι οποίες τέμνονται σε σημείο $\Delta$. Αν $\mathrm{M}$ είναι το μέσον του $\mathrm{A}\mathrm{B}$, να αποδείξετε ότι:
\begin{alist}
    \item Η γωνία $\mathrm{B}\hat{\Delta}\mathrm{A}$ είναι ορθή. \monades{9}
    \item $\mathrm{B}\hat{\mathrm{M}}\Delta = 2 \cdot \mathrm{M}\hat{\Delta}\mathrm{A}$. \monades{8}
    \item $\mathrm{M}\Delta \parallel \varepsilon$. \monades{8}
\end{alist}
